\section{Five Acts Made a Fraternity}

No single date exhausts the foundation. Between 2 July and 18 October 1988, five public instruments moved the project from intention to an erected society with defined liturgical use.

\begin{longtable}{>{\bfseries\RaggedRight\arraybackslash}p{0.17\linewidth}>{\RaggedRight\arraybackslash}p{0.34\linewidth}>{\RaggedRight\arraybackslash}p{0.40\linewidth}}
\toprule
\textbf{Date and act} & \textbf{What it did} & \textbf{What it did not yet do}\\
\midrule
\endfirsthead
\toprule
\textbf{Date and act} & \textbf{What it did} & \textbf{What it did not yet do}\\
\midrule
\endhead
\bottomrule
\endfoot
2 July, declaration of intention & Stated separation from the unauthorized consecrations, fidelity to the pope, and desire for recognized priestly formation and older worship. & It was not an erection decree or final membership roll.\\
18 July, act of foundation at Hauterive & Constituted a projected clerical society under canons 731--746, named it, identified statutory continuity and revision, and petitioned the Holy See. & A private foundation act could not by itself grant pontifical-right status.\\
22 July, PCED declaration & Acknowledged receipt, recognized Bisig's election, and declared the Commission ready to erect the body after review and experimental approval of statutes. & Readiness and recognition of leadership were not the final erection.\\
10 September, liturgical decree & Granted use of the 1962 Missal, Ritual, Pontifical, and Breviary in the body's churches and oratories, with local consent elsewhere except private Masses. & It did not erase local authority or substitute for the fuller government established in October.\\
18 October, Prot. No. 234/88 & Erected the FSSP as a clerical society of apostolic life of pontifical right, approved constitutions for three years, named Bisig for three years, and specified law, mission, liturgy, and communion with bishops. & Experimental constitutions were not the definitive constitutional settlement of 2003.\\
\end{longtable}

The sequence prevents both inflation and minimization. The 22 July letter was more than friendly encouragement: it accepted the foundation record, recognized leadership, and committed the competent commission to a juridical path. Yet the earliest located instrument that says ``erects'' and specifies legal consequences is dated 18 October. The 10 September decree supplied liturgical faculty before that final act; its FSSP-hosted English rendering itself describes the body as of pontifical right, but no earlier declaration of that status was located and the Latin was not independently collated. The study therefore does not use the rendering to move the formal erection date backward.

\subsection{Who counted as a founder?}

The records use different numerical frames. The 2 July declaration contains fifteen names: six styled as priests, Walthard Zimmer styled as a deacon, and eight other signatories whose status the document does not state. The 18 July act displays signatures of eleven priests and Deacon Zimmer; a note says Denis Coiffet, absent that day, is nevertheless considered a founding member, while another says J.-M. Gervais left the proposed foundation before the October erection. The 22 July Roman declaration speaks, in its published English translation, of twelve former SSPX priests.

These formulas need not be forced into one count. One list records a declaration across three cities, another those physically signing a juridical act, another Roman notice, and another institutional memory of the founding cohort. A precise account says which document and moment it means rather than announcing an unexplained canonical number of founders.

\subsection{The name and the promise}

Saint Peter named a program of communion. The 2 July declaration invoked Peter and Paul; the new society's title and later constitutions made attachment to the successor of Peter part of its identity. At the same time, the foundation act preserved continuity of priestly formation and older liturgical practice. The combination was not a temporary bargain waiting for one side to disappear. It was the institution's governing problem and vocation.

\evidenceframe{Founders' act of 18 July; PCED declarations of 22 July and 10 September; erection decree of 18 October 1988.}{Together the instruments identify the society's intended form, Roman acceptance, liturgical faculty, juridical status, initial government, and local-episcopal boundaries.}{The English texts are published by the FSSP and identify German, French, or Latin originals. The exact translations have not been independently collated against authenticated originals; the publication paraphrases rather than reproduces them.}
